\chapter*{About this Document}
\addcontentsline{toc}{chapter}{About this Document}

This document serves as the reference manual for the Libre-V System-on-Chip (SoC). It describes the SoC's functional blocks, their configuration options, and their intended usage. The manual is intended for software developers and hardware designers working with Libre-V platform at any abstraction.

\subsection*{Release information}
\addcontentsline{toc}{section}{Release Information}
\begin{tabularx}{\textwidth}{@{}l l X@{}}
\toprule
\textbf{Version} & \textbf{Date} & \textbf{Changes} \\
\midrule
1.0 & Aug 2026 & Initial release. \\
\bottomrule
\end{tabularx}

\subsection*{Conventions used in this manual}
\addcontentsline{toc}{section}{Conventions}
\begin{itemize}[leftmargin=*]
  \item \textbf{Monospace} text indicates register names, field names, signal names, or code.
  \item \textit{Italic} text indicates a value that you must supply.
  \item Bit ranges are written \texttt{[hi:lo]}, for example \texttt{[7:4]}.
  \item \colorbox{docLightBlue}{\textcolor{docBlue}{\textbf{Blue}}} panels highlight notes; \colorbox{red!5}{\textcolor{docWarnRed}{\textbf{red}}} panels highlight warnings and cautions the user should be careful about.
\end{itemize}

\subsection*{Abbreviations}
\addcontentsline{toc}{section}{Abbreviations}
\begin{tabular}{c l}
  RF & Register File \\
  RO & Read-Only Access \\
  RW & Read-Write Access \\
  WO & Write-Only Access (Reads are undefined/ have no meaning)\\
\end{tabular}

\subsection*{Notes}
\addcontentsline{toc}{section}{Notes}
\begin{itemize}
  \item The official ARM/AMBA specifications have replaced the master--slave terminology with manager--subordinate. However, some parts of the Libre-V were developed before this terminology change and the legacy master--slave nomenclature may appear in corresponding source code.
  \item The base implementation/specification refers to the SoC configuration that you get on cloning the repository. It is possible to configure the SoC to your needs using a Terminal User Interface. Note that this feature is still in beta testing stage and is not recommended for use without caution. 
\end{itemize}
